\medskip

\begin{minipage}{0.6\linewidth}
Les fonctions $f$ et $g$ définies sur $\R$ respectivement par :

\[f(x) =\dfrac12 x^2 +1 \text{ et } g(x) = x + 5,\]

sont représentées sur le graphique ci-contre par la courbe $\mathcal{C}$, courbe représentative de la fonction $f$ et la droite $\Delta$, courbe représentative de la fonction $g$.

Lire graphiquement les positions relatives des courbes représentatives $\mathcal{C}$ de la fonction $f$ et $\Delta$ de la fonction $g$ puis montrer que l'aire de la partie colorée comprise entre la courbe $\mathcal{C}$ et la droite $\Delta$ vaut 18 unités d'aire.
\end{minipage}\hfill
\begin{minipage}{0.39\linewidth}
\psset{unit=0.5cm,arrowsize=2pt 3}
\begin{pspicture*}(-6,-1)(5.5,10)
\pscustom[fillstyle=solid,fillcolor=lightgray]{\psplot[plotpoints=2000,linewidth=1.25pt,linecolor=red]{-2}{4}{x dup mul 0.5 mul 1 add}\psline(4,9)(-2,3)
}
\psgrid[gridlabels=0pt,subgriddiv=1,gridwidth=0.25pt]
\psaxes[linewidth=1.25pt,labelFontSize=\scriptstyle]{->}(0,0)(-6,-1)(5.5,10)
\psplot[plotpoints=2000,linewidth=1.25pt,linecolor=red]{-6}{5.5}{x dup mul 0.5 mul 1 add}
\psplot[plotpoints=500,linewidth=1.25pt,linecolor=blue]{-6}{5.5}{x 5 add}
\uput[r](-4,9){\red $\mathcal{C}$}\uput[ul](2,7){\blue $\Delta$}
\end{pspicture*}
\end{minipage}
